% Generated draft: requires mathematical and source review.
\documentclass[11pt]{article}
\usepackage[margin=25mm]{geometry}
\usepackage{fontspec}
\setmainfont{Latin Modern Roman}
\usepackage{amsmath,amssymb,mathtools,unicode-math}
\setmathfont{Latin Modern Math}
\usepackage{enumitem,needspace,xurl,hyperref}
\setlength{\parindent}{0pt}
\setlength{\parskip}{5pt}
\setlength{\emergencystretch}{3em}
\newcommand{\N}{\mathbb{N}}
\newcommand{\Z}{\mathbb{Z}}
\newcommand{\Q}{\mathbb{Q}}
\newcommand{\R}{\mathbb{R}}
\newcommand{\C}{\mathbb{C}}
\newcommand{\F}{\mathbb{F}}
\newcommand{\A}{\mathbb{A}}
\newcommand{\PP}{\mathbb{P}}
\newcommand{\E}{\mathbb{E}}
\newcommand{\eps}{\varepsilon}
\newcommand{\dd}{\,\mathrm{d}}
\begin{document}
\section*{UnsolvedMath ID: 700368}
\textbf{Source problem code:} catalogue entry
\subsection{Definitions and mathematical statement}
Let \(p\) be a prime. For a finite group \(G\), write \(1_G\) for its identity and \(\operatorname{Aut}(G)\) for its automorphism group. An automorphism \(\alpha\in\operatorname{Aut}(G)\) is fixed-point-free if
\[
\{g\in G:\alpha(g)=g\}=\{1_G\}.
\]
Define the lower central series by \(\gamma_1(G)=G\) and \(\gamma_{i+1}(G)=[\gamma_i(G),G]\), where \([H,G]\) is the subgroup generated by \(h^{-1}g^{-1}hg\) with \(h\in H\) and \(g\in G\). A group is nilpotent of class at most \(c\) if \(\gamma_{c+1}(G)=\{1_G\}\).

The supplied source attributes to G. Higman the assertion that, for every prime \(p\), there is a positive integer bounding the nilpotency classes of all finite groups admitting a fixed-point-free automorphism of order exactly \(p\). Accordingly, define the optimal bound by
\[
\chi_{\mathrm{opt}}(p)=\min\Bigl\{c\in\Z_{\geq1}:\ \forall\text{ finite groups }G,\ \forall\alpha\in\operatorname{Aut}(G),\quad
\bigl(\operatorname{ord}(\alpha)=p\ \text{and}\ \{g\in G:\alpha(g)=g\}=\{1_G\}\bigr)
\Longrightarrow\gamma_{c+1}(G)=\{1_G\}\Bigr\}.
\]
A Higman's function is any function \(\chi\) from the primes to the positive integers satisfying this uniform bound. The source also attributes to Higman the lower bound
\[
\chi(p)\geq\frac{p^2-1}{4}
\qquad\text{for every prime }p
\]
for every such function.

\textbf{Problem.} Determine \(\chi_{\mathrm{opt}}(p)\) for every prime \(p\). In particular, is
\[
\chi_{\mathrm{opt}}(p)=
\begin{cases}
1,&p=2,\\
\dfrac{p^2-1}{4},&p>2
\end{cases}
\]
for every prime \(p\)? The supplied record states that these equalities are known for \(p\leq7\); this historical qualification and the attributed background results are reproduced as supplied, without independent verification.
\subsection{Short English statement}
Determine the smallest uniform bound on the nilpotency class of a finite group admitting a fixed-point-free automorphism of a given prime order \(p\). The proposed optimal bound is \(1\) for \(p=2\) and \((p^2-1)/4\) for odd primes. The supplied source states that this holds for \(p\leq7\).
\subsection{Sources}
\begin{itemize}
\item Local UnsolvedMath registry, supplied complete source record: permanent ID \(700368\), exact source code \texttt{catalogue entry}, title \emph{Kourovka Notebook Problem 6.21}. The registry labels the problem \texttt{open}; this is a supplied status, not an independently verified assessment.
\item \emph{The Kourovka Notebook (Unsolved Problems in Group Theory)}, original problem 6.21, 6th issue (1978); proposed by V. D. Mazurov. The supplied record identifies the 21st issue edition (2026), arXiv:1401.0300v48: \url{https://arxiv.org/abs/1401.0300v48}. The collection description names E. I. Khukhro and V. D. Mazurov as editors.
\item The background existence assertion and lower bound are attributed to G. Higman in the supplied statement. No separate bibliographic citation for Higman's work is supplied.
\end{itemize}
\end{document}
